\documentclass[12pt]{extarticle}
\def\activityid{F04-FAC-v3}\usepackage[letterpaper,margin=.65in,top=.60in,bottom=.65in]{geometry}
\usepackage{fix-cm}
\usepackage[T1]{fontenc}\usepackage[default]{sourcesanspro}
\usepackage{microtype,tikz,array,tabularx,fancyhdr,amsmath,amssymb}
\usepackage[hidelinks]{hyperref}\usetikzlibrary{calc,arrows.meta}
\setlength{\parindent}{0pt}\setlength{\parskip}{7pt}\setlength{\headheight}{0pt}\setlength{\footskip}{26pt}
\setlength{\emergencystretch}{2em}\pagestyle{fancy}\fancyhf{}\renewcommand{\headrulewidth}{0pt}\renewcommand{\footrulewidth}{.4pt}
\fancyfoot[L]{\scriptsize Bellingham Math Circle / Week 4 / \activityid}\fancyfoot[R]{\scriptsize \thepage}
\definecolor{ink}{gray}{.12}\definecolor{soft}{gray}{.95}\color{ink}
\newcommand{\PageTitle}[2]{{\fontsize{10}{12}\selectfont\bfseries WEEK 4\quad ROUND AND ROUND\hfill #1\par}\vspace{3pt}{\fontsize{26}{29}\selectfont\bfseries #2\par}\vspace{2pt}}
\newcommand{\NameLine}{{\small Name: \rule{2.65in}{.4pt}\hfill Date: \rule{1.35in}{.4pt}\par}\vspace{4pt}}
\newcommand{\Task}[2]{\par\vspace{3pt}{\large\bfseries #1}\par #2\par}
\newcommand{\Subhead}[1]{\par\vspace{3pt}{\large\bfseries #1}\par}
\newcommand{\RuleBox}[1]{\par\begingroup\setlength{\fboxsep}{8pt}\colorbox{soft}{\parbox{\dimexpr\linewidth-16pt\relax}{#1}}\endgroup\par}
\newcommand{\WriteBox}[2]{\par\textbf{#1}\par\vspace{2pt}\begin{tikzpicture}\draw[black!40,line width=.5pt] (0,0) rectangle (\dimexpr\linewidth-.5pt\relax,#2);\end{tikzpicture}\par}
\newcommand{\StudentFooter}{\vfill{\small Try it with objects. Explain by pointing, talking, or drawing.}\par}
\newcommand{\LampRing}[3]{\begin{tikzpicture}[x=1in,y=1in]
\foreach \i in {1,...,#1}{\coordinate (v\i) at ({90-360*(\i-1)/#1}:#2);}
\foreach \i in {1,...,#1}{\pgfmathtruncatemacro{\j}{mod(\i,#1)+1}\draw[thick] (v\i)--node[midway,fill=white,inner sep=2pt,font=\small]{\i\j}(v\j);}
\foreach \i in {1,...,#1}{\node[circle,draw,fill=white,minimum size=.52in,inner sep=0pt,font=\bfseries] at(v\i){\i};}
\foreach \i in {#3}{\node[circle,draw,fill=ink,text=white,minimum size=.52in,inner sep=0pt,font=\bfseries] at(v\i){\i};}
\end{tikzpicture}}
\newcommand{\Lights}[1]{\begin{tikzpicture}[x=.55in,y=.55in,baseline=-.10in]\foreach \v [count=\i] in {#1}{\ifnum\v=1\fill (\i,0) circle(.16in);\else\draw (\i,0) circle(.16in);\fi}\end{tikzpicture}}
\newcommand{\ClockRing}[2]{\begin{tikzpicture}[x=1in,y=1in]\draw[black!30] (0,0) circle(#2);\pgfmathtruncatemacro{\n}{#1-1}\foreach\i in {0,...,\n}{\node[circle,draw,fill=white,minimum size=.35in,inner sep=0pt] at({90-360*\i/#1}:#2){\i};}\draw[-{Stealth},thick] (-.35,.15) arc(155:25:.38);\node[font=\small] at(0,-.18){clockwise};\end{tikzpicture}}
\newcommand{\Key}[2]{\begin{center}\renewcommand{\arraystretch}{1.55}\begin{tabular}{c|*{#1}{c}}in&#2\end{tabular}\end{center}}


% v3 student pages use one running header and numbered tasks only.
\newcommand{\StudentHeader}[1]{{\fontsize{10}{12}\selectfont\bfseries Week 4 / Ring shifts / #1\par}\vspace{10pt}}
\newcommand{\Problem}[2]{\par\vspace{4pt}\textbf{Problem #1:} #2\par}
\newcommand{\Space}[1]{\begin{tikzpicture}\draw[black!35] (0,0) rectangle (\dimexpr\linewidth-.5pt\relax,#1);\end{tikzpicture}\par}
\newcommand{\Record}[2]{#1\par\Space{#2}}

\begin{document}

\PageTitle{FACILITATOR / 1}{Prepare the ring shifts}
\RuleBox{v3 revision, September 27, 2026: unpiloted. This applies Week 1's classroom guidance to new ring and shuffle work; no outcomes from teaching this revision are known.}
\Subhead{Materials and prerequisites}
Print first pages for three K, three middle, and one upper child (seven starting sheets). Keep continuation masters ready to copy as needed, including the later ring mats. Leave the extra in reserve. Bring seven counters, pencils, scratch paper, and cards 0--7; cards 8--15 extend the reserve. A paper strip numbered 0--7 distinguishes card places from labels. No wheels or brads are required.\par
\textbf{K:} match four shapes, count hops to three with support. \textbf{Middle:} labels A--J, count to ten; read-aloud works. \textbf{Upper:} count to eighteen and compare routes; general gcd proof additionally requires division and adult-supported coprime divisibility. \textbf{Extra:} track positions and repeated doubling with remainders; exponent notation is optional.
\Subhead{Handle, then launch: 0--15 minutes}
Give 0--10 minutes to explore counters on rings. At 10--15 minutes, gather around the four-shape mat. Define the arrow direction by moving circle to triangle. Show a two-hop turn, counting only movements, and point to square as the landing; triangle was passed over. Reset and let a child replay it. Draw only the landing after the starting circle. Take one more two-hop turn and record its landing. Now distribute different starting tasks. The purpose is a shared action and record, not a gcd explanation.
\Subhead{A flexible hour}
\textbf{15--35:} K compares starts under one/two hops; middle uses the common shift and records two routes; upper explores +4 and +8. Adult A anchors K; organizer checks landing records with older children. Compare routes when useful. \textbf{35--40:} stand, stretch, reset. \textbf{40--55:} continue experiments or introduce one new record when useful. \textbf{55--60:} share and tidy. Page counts are reserves, not a completion schedule.
\Subhead{Stops and later directions}
\textbf{K:} stop after one-hop and two-hop comparisons; three hops and hidden-turn undo are later choices. \textbf{Middle:} stop after +2 routes from A,B,C,D. Marking the two sets helps explain a split already encountered; +3,+5 and inverse checks follow. The optional complete list of full-visit steps is an adult extension.\par
\textbf{Upper:} +4/+8 loops are a satisfying stop. Compare +3,+6,+5 before the summary table (7). Skip filling known entries if copying slows the child. Offer distance records only after physical routes; the strip/gcd proof is optional adult work. \textbf{Extra:} obtain several real rows and position tracks before the map, loops, or doubling rule. Stop at the eight-card return if that is enough.
\Subhead{What to listen for}
``Did you land there or pass over it?'' ``Which start is unused?'' ``Is this a card label or its place?'' All shifts can be undone; a short orbit does not mean an undecodable key. Discuss why an argument covers every start only after children can show examples. Proofs and checked outputs follow.
\newpage
\PageTitle{FACILITATOR / 2}{Exact routes and the gcd rule}
\Subhead{K--1}
From circle, one hop visits triangle, square, star, circle (four turns); two hops alternate circle/square (two turns), never landing on triangle or star. A secret circle-to-square shift is two hops, so triangle becomes star, square becomes circle, and star becomes triangle. Two hops undo themselves. Three-hop turns follow circle, star, square, triangle, circle, visiting all four. On six positions, among steps 1 through 5, only 1 and 5 visit all six; 2 and 4 give two loops of length 3; 3 gives three loops of length 2.
\Subhead{Middle}
A to D is +3 on A--J. GDG decodes to \textbf{DAD}; BAG encodes to \textbf{EDJ}. B becomes E, so the additional B-to-F clue contradicts the fixed-shift rule. Repeated +2 gives A,C,E,G,I,A and B,D,F,H,J,B. Marking alternate letters explains why the two routes cannot meet: adding two preserves that marking. Repeated +3 gives A,D,G,J,C,F,I,B,E,H,A. The steps visiting all ten are \textbf{1,3,7,9}. Any shift is reversible: +2 is undone by -2, or +8. Skipping some letters along one repeated orbit is not information loss.
\Subhead{Upper: key and complete classification}
9 to 1 gives +4, which takes 2 to 6. Outputs 1,4,10 decode to \textbf{9,0,6}. Its four loops are (0,4,8), (1,5,9), (2,6,10), (3,7,11), each returning in three turns.
\begin{center}\renewcommand{\arraystretch}{1.25}\begin{tabular}{c|rrrrrrrrrrrr}Shift k&0&1&2&3&4&5&6&7&8&9&10&11\\\hline Return time&1&12&6&4&3&12&2&12&3&4&6&12\\Number of loops&12&1&2&3&4&1&6&1&4&3&2&1\end{tabular}\end{center}
Steps 1,5,7,11 visit all twelve. The zero rule is the identity and has first \emph{positive} return time 1, not zero. For eighteen positions, +6 returns after 3; +5 after 18.
\Subhead{Why the formula works}
For +8 on 12 places, return means \(3\mid2t\). Draw three t-strips, totaling 3t. If the first two strips total a multiple of 3, removing them leaves t a multiple of 3. Thus the first positive return is 3. This proves this case for every t. The optional general bridge is on page 4; do not treat a completed table as its proof.
\newpage
\PageTitle{FACILITATOR / 3}{From turning to shuffling}
\Subhead{Upper: inverse and composition}
On twelve places, +4 is undone by +8 (or -4). For every starting x, +4 followed by +7 gives x+11 modulo12; reversing the order gives the same because integer addition commutes. Every shift is bijective, even if its orbit from 0 is a proper subset. The visit-everywhere condition is \(\gcd(n,k)=1\), not a condition for reversibility. This distinction is central: a substitution permutation can have several cycles.
\Subhead{Extra: eight cards, three shuffles}
Rows, left to right with the leftmost card the top:
\begin{center}\renewcommand{\arraystretch}{1.5}\begin{tabular}{c|l}Shuffle&Row\\\hline 0&0 1 2 3 4 5 6 7\\1&0 4 1 5 2 6 3 7\\2&0 2 4 6 1 3 5 7\\3&0 1 2 3 4 5 6 7\end{tabular}\end{center}
The old-to-new position map is \textbf{0,2,4,6,1,3,5,7}. Its cycles are fixed 0, fixed 7, (1,2,4), and (3,6,5). Thus the whole shuffle returns after 3 and not earlier. This is also a direct week 3 loop argument.
For old position x in the first half (0--3), the new position is 2x. In the second half (4--7), it is \(2(x-4)+1=2x-7\). For x below 7 these are exactly \(2x\bmod7\). Position 7 stays 7, separately; it must not be merged with residue 0. Repetition gives \(2^t x\bmod7\). The powers 2,4,8 have remainders 2,4,1, so the first identity occurs at t=3. Tracking position 1 rules out earlier global returns, and remainder 1 fixes all positions 0--6.
\Subhead{Extra predictions}
Six cards use modulus 5. Powers 2,4,8,16 give remainders 2,4,3,1, so the order is \textbf{4}. Sixteen cards use modulus 15: 2,4,8,16 give 2,4,8,1, again order \textbf{4}. In general an out-shuffle on 2m cards maps positions below \(2m-1\) by doubling modulo \(2m-1\), fixes the final position, and has the multiplicative order of 2 as return time (the two-card identity is a trivial edge case). More cards can therefore have the same or even a smaller return time.
\Subhead{A useful additional contrast}
Random shuffling asks about a distribution over possible orders. This perfect shuffle is one precisely defined permutation repeated deterministically. Its return time says nothing by itself about how mixed a random shuffle is. Do not introduce chance unless it becomes a separate, explicitly changed rule.
\newpage
\PageTitle{FACILITATOR / 4}{Mathematical lineage and reuse}
\Subhead{Optional adult bridge: coprime divisibility}
If positive a,b are coprime and \(a\mid bt\), compare base strips a,b with scaled strips at,bt. Repeatedly subtract the smaller base strip from the larger, doing the same to the scaled strips. Both scaled lengths stay multiples of a. The positive base sum decreases, so the process ends with equal lengths c. Common divisors are unchanged: a divisor of x,y also divides x-y,y, and conversely by adding. Thus c=1, since the original lengths were coprime. The scaled process ends at t, which is therefore a multiple of a. Conversely, \(a\mid t\) plainly implies \(a\mid bt\).
\par
Try 5,3 first: 5-3=2, then 3-2=1. On scaled strips this leaves \(3t-(5t-3t)=t\). Equivalently, \(2\cdot3t-5t=t\); if \(5\mid3t\), the remainder is a multiple of 5. Ask the child to explain each removal. This is additional readiness, not assumed knowledge of prime factors.
\par
For k>0, set \(d=\gcd(n,k)\), \(n=da\), \(k=db\). Return means \(a\mid bt\); the lemma gives \(a\mid t\), hence first return n/d. Every start has this length, so there are d loops. For k=0, every point is fixed. If only the 3-and-2 case was explained, record that proof and label the general rule conjectured or given as a theorem. Full strip prompts and the Week 9 application: \texttt{plans/coprime-divisibility-scaffold.md}.
\Subhead{Sources and boundaries}
\textbf{Judson, \emph{Abstract Algebra: Theory and Applications}}, Ch. 4, \S4.1, Theorem 4.13, pp. 48--49, gives the cyclic order formula: \href{https://people.hsc.edu/faculty-staff/blins/books/JudsonAbstractAlgebra.pdf}{textbook}. Our rings are its additive version.
\textbf{Diaconis, Graham, Kantor (1983)}, \emph{The mathematics of perfect shuffles}, \S2, Lemma 1, p. 177, proves the out-shuffle order via doubling: \href{https://pages.uoregon.edu/kantor/PAPERS/PerfectShuffles.pdf}{original paper}. The extra develops this small return-time case and position rule, not the paper's shuffle-group classification. The plan has full research references.
Rozhkovskaya, Lesson 2, problems 2.3--2.7, supplies the shift-cipher entry. \emph{Math Circle by the Bay}, preface viii--x, supports manipulatives, changing pace, and reserves; Rozhkovskaya Lesson 8 records copying slowing the activity. Our rings, proof scaffolds, and hour are adaptations. Lowell Handout 6, problem 6.7, and Handout 7, problem 7.8, use a different card-ordering solitaire. Ask about recognition; do not assume this shuffle was taught.
\Subhead{Use history and checks}
Record actual steps, ring sizes, and extensions used as F04-K/M/U/X-v3. Keep 26-letter Caesar messages as a later alternative, not extra arithmetic for early finishers. Future years can compare prime/composite ring sizes, affine maps, in-shuffles, or perfect-shuffle groups; avoid replaying the same small cases.
\texttt{verify.py} checks all 12 shifts, the 10-letter routes, message examples, the 18-ring predictions, shuffle rows, position formulas, and first return times for 6, 8, 16 cards. The facilitator proofs establish why the patterns hold.

\newpage
\PageTitle{FACILITATOR / 5}{v3 checks and observations}
\Subhead{Added ring experiments}
\textbf{K 1--2:} from circle land triangle,square,star,circle; from triangle land square,star,circle,triangle. \textbf{3--4:} circle and square share the two-hop orbit; triangle and star share the other. Changing from circle to square never reaches triangle. \textbf{5--7:} the three-hop, hidden shift, and inverse results are on page 2. Read and scribe the shape records; writing is not a gate.\par
\textbf{Middle 1:} +2 takes A to C and I to A. \textbf{2:} +3 takes B,I,J to E,B,C. \textbf{3--4:} messages and incompatible clue are on page 2. \textbf{5--7:} C follows C,E,G,I,A,C, the A-route; D follows D,F,H,J,B,D, the B-route. The two markings are preserved by every +2 turn. \textbf{8:} full +3 route is on page 2. \textbf{9:} +5 gives A,F,A and B,G,B, length 2. \textbf{10:} +2 then +8 gives A,C,A; B,D,B; I,A,I; J,B,J. Each starting letter returns despite the shorter orbits.
\Subhead{Upper: new comparisons}
\textbf{1:} +4 sends 0,9,2,8 to 4,1,6,0. \textbf{4:} +8 traverses each +4 loop in reverse. \textbf{5:} +3 gives three loops of length 4, +6 six loops of length 2, +5 one loop of length 12. \textbf{7:} full table is on page 2. \textbf{8:} eighteen-place +6 route is 0,6,12,0; +5 is 0,5,10,15,2,7,12,17,4,9,14,1,6,11,16,3,8,13,0. \textbf{9:} distances at turns 1--4 are 4,8,12,16 and 8,16,24,32; first whole laps are turn 3 in both. \textbf{10:} +4 then +8 returns all starts, with intermediate positions 4,5,9,3. \textbf{11:} both orders total +11, sending 0,2,9 to 11,1,8. \textbf{12:} in general shifts +a,+b compose as +(a+b), wrapping modulo 12.
\Subhead{Extra: position scaffolding}
\textbf{1--6:} page 3 gives all eight-card rows, map, loops, doubling proof, and fixed last position. Position traces from starts 1,3,7 are 1,2,4,1; 3,6,5,3; and always 7. \textbf{7:} six-card rows are 012345,031425,043215,024135,012345. \textbf{8:} sixteen-card loops are (0),(15),(1,2,4,8),(3,6,12,9),(5,10),(7,14,13,11). The full return is 4. \textbf{9:} 6 and 16 cards both take 4; 8 cards take only 3. Ask for a proof after these checks: position 1 rules out earlier returns, and the same remainder-1 multiplier fixes all other positions below the last.
\Subhead{Evidence for the next revision}
Record actual step sizes/decks attempted, mistaken starting-point counts, skipped-versus-landed confusion, success using a route record, and card-label/position confusion. Note exact adult rescue and what children wanted more time for. Distinguish observations from hypotheses, and cases checked from arguments proved. The extra and general theorem may remain unused; this is not incomplete work. v3 is unpiloted.

\end{document}
